
\documentclass{article}
\usepackage{colortbl}
\usepackage{makecell}
\usepackage{multirow}
\usepackage{supertabular}

\begin{document}

\newcounter{utterance}

\centering \large Interaction Transcript for game `cladder', experiment `full\_v1.5\_default', episode 7668 with selfrefine{-}llama{-}3.1{-}8b{-}instruct.
\vspace{24pt}

{ \footnotesize  \setcounter{utterance}{1}
\setlength{\tabcolsep}{0pt}
\begin{supertabular}{c@{$\;$}|p{.15\linewidth}@{}p{.15\linewidth}p{.15\linewidth}p{.15\linewidth}p{.15\linewidth}p{.15\linewidth}}
    \# & \multicolumn{2}{c}{Player} && \multicolumn{2}{c}{Game Master} \\
    \hline

    \theutterance \stepcounter{utterance}  
    & & & \multicolumn{4}{p{0.6\linewidth}}{
        \cellcolor[rgb]{0.9,0.9,0.9}{
            \makecell[{{p{\linewidth}}}]{
                \texttt{\tiny{[P1$\langle$GM]}}
                \texttt{You are an expert in causal inference. The following question is not a typical commonsense query, but rather a meticulously designed question created by a professor specializing in causal inference, intended to assess the students' mastery of the course content.} \\
\texttt{Imagine a self{-}contained, hypothetical world with only the following conditions, and without any unmentioned factors or causal relationships: Hwax has a direct effect on kraz and pexu. Pexu has a direct effect on rukz. Kraz has a direct effect on rukz. The overall probability of pexu is 41\%. For those who are not pexu, the probability of rukz is 66\%. For those who are pexu, the probability of rukz is 22\%. Is rukz more likely than not rukz overall?} \\
\texttt{Start your answer with 'Yes' or 'No', followed by additional reasoning or evidence to support your explanation.} \\
            }
        }
    }
     \\ \\

    \theutterance \stepcounter{utterance}  
    & & \multicolumn{4}{p{0.6\linewidth}}{
        \cellcolor[rgb]{0.95,0.95,0.95}{
            \makecell[{{p{\linewidth}}}]{
                \texttt{\tiny{[P1$\leftrightarrow$P]}}
                \texttt{{[}DRAFT{]}} \\
\texttt{No, rukz is not more likely than not rukz overall.} \\
\\ 
\texttt{To determine the overall probability of rukz, we need to consider the probabilities of rukz given different values of pexu and the probability of pexu itself. } \\
\\ 
\texttt{Let's denote the probability of rukz given pexu as P(Rukz|Pexu) and the probability of rukz given not pexu as P(Rukz|\textasciitilde{}Pexu). We are given that P(Rukz|\textasciitilde{}Pexu) = 66\% and P(Rukz|Pexu) = 22\%.} \\
\\ 
\texttt{We also know that P(Pexu) = 41\%. To find the probability of not pexu, we can subtract this from 1: P(\textasciitilde{}Pexu) = 1 {-} P(Pexu) = 1 {-} 0.41 = 0.59.} \\
\\ 
\texttt{Now, we can use the law of total probability to find the overall probability of rukz:} \\
\\ 
\texttt{P(Rukz) = P(Rukz|Pexu) * P(Pexu) + P(Rukz|\textasciitilde{}Pexu) * P(\textasciitilde{}Pexu)} \\
\texttt{= 0.22 * 0.41 + 0.66 * 0.59} \\
\texttt{= 0.0902 + 0.3894} \\
\texttt{= 0.4796} \\
\\ 
\texttt{Since 0.4796 is less than 0.5, rukz is not more likely than not rukz overall.} \\
            }
        }
    }
    & \\ \\

    \theutterance \stepcounter{utterance}  
    & & \multicolumn{4}{p{0.6\linewidth}}{
        \cellcolor[rgb]{0.95,0.95,0.95}{
            \makecell[{{p{\linewidth}}}]{
                \texttt{\tiny{[P1$\leftrightarrow$P]}}
                \texttt{{[}SELF{-}FEEDBACK{]}} \\
\texttt{FORMAT: ok} \\
\texttt{RULES: ok} \\
\texttt{STRATEGY: ok} \\
\texttt{REASON: The answer correctly applies the law of total probability to calculate the overall probability of rukz and provides a clear explanation of the reasoning.} \\
            }
        }
    }
    & \\ \\

    \theutterance \stepcounter{utterance}  
    & \multicolumn{4}{p{0.6\linewidth}}{
        \cellcolor[rgb]{0.561,0.737,0.561}{
            \makecell[{{p{\linewidth}}}]{
                \texttt{\tiny{[P1$\rangle$GM]}}
                \texttt{ok} \\
\texttt{RULES: ok} \\
\texttt{STRATEGY: ok} \\
\texttt{REASON: The answer correctly applies the law of total probability to calculate the overall probability of rukz and provides a clear explanation of the reasoning.} \\
\\ 
\texttt{No, rukz is not more likely than not rukz overall.} \\
\\ 
\texttt{To determine the overall probability of rukz, we need to consider the probabilities of rukz given different values of pexu and the probability of pexu itself. } \\
\\ 
\texttt{Let's denote the probability of rukz given pexu as P(Rukz|Pexu) and the probability of rukz given not pexu as P(Rukz|\textasciitilde{}Pexu). We are given that P(Rukz|\textasciitilde{}Pexu) = 66\% and P(Rukz|Pexu) = 22\%.} \\
\\ 
\texttt{We also know that P(Pexu) = 41\%. To find the probability of not pexu, we can subtract this from 1: P(\textasciitilde{}Pexu) = 1 {-} P(Pexu) = 1 {-} 0.41 = 0.59.} \\
\\ 
\texttt{Now, we can use the law of total probability to find the overall probability of rukz:} \\
\\ 
\texttt{P(Rukz) = P(Rukz|Pexu) * P(Pexu) + P(Rukz|\textasciitilde{}Pexu) * P(\textasciitilde{}Pexu)} \\
\texttt{= 0.22 * 0.41 + 0.66 * 0.59} \\
\texttt{= 0.0902 + 0.3894} \\
\texttt{= 0.4796} \\
\\ 
\texttt{Since 0.4796 is less than 0.5, rukz is not more likely than not rukz overall.} \\
            }
        }
    }
    & & \\ \\

    \theutterance \stepcounter{utterance}  
    & & & \multicolumn{2}{p{0.3\linewidth}}{
        \cellcolor[rgb]{0.9,0.9,0.9}{
            \makecell[{{p{\linewidth}}}]{
                \texttt{\tiny{[GM$|$GM]}}
                \texttt{ParseError: The response doesn't start with one of the required word {[}'yes', 'no'{]}, but ok} \\
            }
        }
    }
    & & \\ \\

    \theutterance \stepcounter{utterance}  
    & & & \multicolumn{2}{p{0.3\linewidth}}{
        \cellcolor[rgb]{0.9,0.9,0.9}{
            \makecell[{{p{\linewidth}}}]{
                \texttt{\tiny{[GM$|$GM]}}
                \texttt{game\_result = ABORT} \\
            }
        }
    }
    & & \\ \\

\end{supertabular}
}

\end{document}
